% !TeX root = ../../main.tex
\subsection{\texorpdfstring{$t$-based inference}{t-based inference}}
\label{sec:inference}

At convergence, the model-based covariance estimator is
\begin{equation}
  \widehat{\operatorname{Cov}}
  (\widehat{\boldsymbol{\beta}})
  =(X^\top\widehat{W}X)^{-1}.
  \label{eq:covariance}
\end{equation}
The Pearson equation scales the residual variation to approximately one.  If
the dispersion estimate reaches its numerical upper boundary, the implementation multiplies Equation~\eqref{eq:covariance} by the
remaining Pearson scale.

The plug-in covariance in Equation~\eqref{eq:covariance} treats the estimated
$\widehat{\rho}$ as fixed and omits uncertainty from estimating a separate
dispersion for every guide.  Referring the coefficient to $t_{m-q}$ gives a
heavier-tailed small-sample reference than a normal approximation, but it is
not a formal correction for dispersion-estimation uncertainty.  With few
independent libraries, calibration must therefore be checked by simulation,
phenotype permutation, or prespecified negative controls.

For coefficient $\beta_j$, the statistic
\begin{equation}
  t_j =
  \frac{\widehat{\beta}_j-\beta_{j,0}}
  {\sqrt{
  \widehat{\operatorname{Cov}}
  (\widehat{\boldsymbol{\beta}})_{jj}}}
  \label{eq:tstat}
\end{equation}
is compared with a Student $t_{m-q}$ distribution.  The number of independent sequenced samples determines the degrees of freedom;
total read count does not enter that calculation.  Together with the bounded weights, this prevents sequencing depth
from being treated as biological replication.

More generally, for a prespecified contrast vector $\mathbf{c}$,
\begin{equation}
  t_{\mathbf{c}} =
  \frac{\mathbf{c}^\top\widehat{\boldsymbol{\beta}}-\delta_0}
  {\sqrt{\mathbf{c}^\top
  \widehat{\operatorname{Cov}}(\widehat{\boldsymbol{\beta}})
  \mathbf{c}}}
\end{equation}
uses the same $m-q$ degrees of freedom.  This covers pairwise group contrasts,
average slopes in interaction models, and other one-degree-of-freedom
hypotheses.  Guide-wise two-sided $p$-values are adjusted using the
Benjamini--Hochberg procedure \citep{benjamini1995controlling}.

\subsubsection{Optional guide-dispersion moderation}

BARCS-MOD applies an optional empirical-Bayes transformation to the fitted
guide statistics while retaining the same likelihood.  Let $\nu=m-q$ and let
$s_g^2=Q_g^{(0)}/\nu$ be guide $g$'s Pearson variance inflation under the
binomial working model.  A lowess trend in log mean CPM gives a guide-specific
prior scale $s_{0g}^2$, and the scaled-$F$ log-moment estimator of
\citet{smyth2004linear} supplies prior degrees of freedom $d_0$.  The moderated
inflation is
\begin{equation}
  \widetilde{s}_g^2
  =
  \frac{\nu s_g^2+d_0s_{0g}^2}{\nu+d_0}.
  \label{eq:guide-moderation}
\end{equation}
Because the original beta-binomial fit cannot estimate negative
overdispersion, its stored standard error corresponds to the fitted inflation
$s_{g,\mathrm{fit}}^2=\max(1,s_g^2)$.  The implementation therefore reports
\begin{equation}
  \operatorname{SE}_{g,\mathrm{mod}}
  =
  \operatorname{SE}_{g}
  \sqrt{\frac{\widetilde{s}_g^2}{\max(1,s_g^2)}},
  \qquad
  t_{g,\mathrm{mod}}
  =
  \frac{\widehat\beta_g}{\operatorname{SE}_{g,\mathrm{mod}}},
  \qquad
  \nu_{\mathrm{mod}}=\nu+d_0 .
  \label{eq:guide-moderated-t}
\end{equation}
The \texttt{pmax} truncation in
Equation~\eqref{eq:guide-moderated-t} retains the fitted inflation floor in
the standard-error rescaling.  The optional \texttt{one\_way=TRUE}
variant replaces $\widetilde{s}_g^2$ by
$\max(\widetilde{s}_g^2,s_{g,\mathrm{fit}}^2)$, so moderation can only increase
a standard error.  All reported BARCS-MOD results are simulation ablations;
the real-data benchmarks use unmoderated BARCS.

\subsubsection{From guide statistics to gene statistics}

BARCS fits the count model guide by guide.  The Cas13, FACS, CRISPulator, GSE70038,
A375, and descriptive HT-29 benchmark scripts use an explicit
directional Stouffer summary, denoted BARCS-ST below, to obtain one gene-level
statistic.  For the $j$th valid guide targeting gene $g$, the two-sided guide
$p$-value is converted back to a signed standard-normal score,
\begin{equation}
  z_{gj}
  =
  \operatorname{sign}(\widehat{\beta}_{gj})
  \Phi^{-1}\!\left(1-\frac{p_{gj}}{2}\right).
  \label{eq:signedz}
\end{equation}
If $m_g$ valid guides target the gene, their combined score and two-sided
gene-level $p$-value are
\begin{equation}
  Z_g=\frac{1}{\sqrt{m_g}}\sum_{j=1}^{m_g}z_{gj},
  \qquad
  p_g=2\Phi(-|Z_g|).
  \label{eq:genestouffer}
\end{equation}
The reported gene effect is the median of the guide coefficient estimates,
\begin{equation}
  \widehat{\beta}_g
  =\operatorname{median}_{j=1,\ldots,m_g}
    \widehat{\beta}_{gj},
  \label{eq:geneeffect}
\end{equation}
and gene-level false-discovery rates are obtained by applying
Benjamini--Hochberg to the collection of $p_g$ values.  Concordant guide directions reinforce one another in the combined score, and
discordant directions cancel.  The median limits the influence of one extreme
guide on the reported effect.

This summary does not fit a hierarchical gene model.
Equations~\eqref{eq:signedz}--\eqref{eq:genestouffer} give equal weight to
valid guides and use the independence reference $\sqrt{m_g}$; shared sequence
artifacts or other correlation can therefore make the gene statistic
anti-conservative.  An extension for gene-level inference should estimate
guide reliability and within-gene dependence or use a hierarchical effect
model.  In the Cas13 comparison, the four guide-level regression methods use
signed Stouffer aggregation as specified below.  In the other biological
head-to-head benchmarks, the Stouffer--median procedure is used for BARCS (and
for the naive binomial comparator).  MAGeCK, Chronos, BAGEL2, Waterbear, and
MAUDE retain their native or deposited gene-level summaries.  The exploratory
DeWeirdt false-discovery audit used Fisher aggregation as stated in its
Results description and is not pooled with these primary benchmark statistics.
